\section{Recursive queries}\label{sec:recursion}

Recursive queries are very useful in a many applications.
For example, many graph algorithms (such as graph reachability
or transitive closure) are naturally expressed using recursive queries.

We introduce two new stream operators that are instrumental in
expressing recursive query evaluation.  These operators allow us
to build circuits implementing looping constructs, which 
are used to iterate computations until a fixed-point is reached.

\begin{definition}\label{def:zae}
We say that a stream $s \in \stream{A}$ is \defined{zero almost-everywhere} if it has a finite 
number of non-zero values, i.e., there exists a time $t_0 \in \N$
s.t. $\forall t \geq t_0 . s[t] = 0$.
\noindent Denote the set of streams that are zero almost everywhere
by $\streamf{A}$.
\end{definition}

\paragraph{Stream introduction}

The delta function (named from the Dirac delta function) $\delta_0 : A \rightarrow \stream{A}$
produces a stream from a scalar value:
$$\delta_0(v)[t] \defn \left\{
\begin{array}{ll}
  v & \mbox{if } t = 0 \\
  0_A & \mbox{ otherwise}
\end{array}
\right.
$$
\ifstreamexamples
For example, $\delta_0(5)$ is the stream $\sv{5 0 0 0 0}$.
\fi

\begin{comment}
Here is a diagram showing a $\delta_0$ operator; note that the input is a scalar value,
while the output is a stream:  

\begin{tikzpicture}[auto,node distance=1cm,>=latex]
    \node[] (input) {$i$};
    \node[block, right of=input] (delta) {$\delta_0$};
    \node[right of=delta] (output) {$o$};
    \draw[->] (input) -- (delta);
    \draw[->] (delta) -- (output);
\end{tikzpicture}
\end{comment}

\paragraph{Stream elimination}

We define the function $\int : \streamf{A} \rightarrow
A$, over streams that are zero almost everywhere, as 
$\int(s) \defn \sum_{t \geq 0} s[t]$.  
$\int$ is closely related to $\I$; if $\I$ is the
indefinite integral, $\int$ is the definite integral on the
interval $0 - \infty$.

For many implementation strategies (including  
relational and Datalog queries given below) the $\int$
operator can be approximated finitely and precisely by integrating until the
first 0 value encountered, since it can be proven that its input is always
the derivative of a monotone stream.

\begin{comment}
Here is a diagram using the $\int$ operator; note that  the result it 
produces is a scalar, and not a stream:

\begin{tikzpicture}[auto,node distance=1cm,>=latex]
    \node[] (input) {$i$};
    \node[block, right of=input] (S) {$\int$};
    \node[right of=S] (output) {$o$};
    \draw[->] (input) -- (S);
    \draw[->] (S) -- (output);
\end{tikzpicture}
\end{comment}

$\delta_0$ is the left inverse of $\int$, i.e.: $\int \circ \; \delta_0 = \id_A$.  
\begin{proposition}
$\delta_0$ and $\int$ are LTI.
\end{proposition}

\paragraph{Nested time domains}

So far we used a tacit assumption that ``time'' is common for all
streams in a program.  For example, when we add two streams, 
we assume that they use the same ``clock'' for the time dimension.
However, the $\delta_0$ operator creates a stream with a ``new'', independent time
dimension.  We require \emph{well-formed circuits}
to ``insulate'' such
nested time domains by nesting them between a $\delta_0$ and an $\int$ operator:

\begin{center}
\begin{tikzpicture}[auto,node distance=1cm,>=latex]
    \node[] (input) {$i$};
    \node[block, right of=input] (delta) {$\delta_0$};
    \node[block, right of=delta] (f) {$Q$};
    \draw[->] (input) -- (delta);
    \draw[->] (delta) -- (f);

    \node[block, right of=f] (S) {$\int$};
    \node[right of=S] (output) {$o$};
    \draw[->] (f) -- (S);
    \draw[->] (S) -- (output);
\end{tikzpicture}
\end{center}

\begin{proposition}
If $Q$ is time-invariant, the circuit above has the zero-preservation
property: $\zpp{\int \circ\; Q \circ \delta_o}$.
\end{proposition}

\subsection{Implementing Recursive Datalog}\label{sec:datalog}

We illustrate the implementation of recursive queries in \dbsp for
stratified Datalog.
Datalog is strictly more expressive than
relational algebra since it can express recursive programs, e.g.:
{\small
\begin{lstlisting}[language=ddlog]
O(v) :- I(v).  // base case
O(v) :- I(z), O(x), v = ... .  // rec case
\end{lstlisting}
}
In general, a recursive Datalog program defines a set of
mutually recursive relations $O_1,..,O_n$ as an equation
$(O_1,..,O_n)=R(I_1,..,I_m, O_1,..,O_n)$, where $I_1,..,I_m$ are
input relations and $R$ is a relational (non-recursive) query.

The following algorithm generates a circuit that computes a
solution to this equation.  We describe the algorithm informally and 
for the special case of a single input $I$ and single output $O$; the general case
can be found in the companion technical report~\cite{tr}, and is only 
slightly more involved.

\noindent 
\begin{enumerate}[nosep, leftmargin=\parindent]
\item Implement the non-recursive relational query $R$ as described in
    \secref{sec:relational} and Table~\ref{tab:relational}; this produces
    an acyclic circuit whose inputs and outputs are a \zr (i.e., not a stream):
    \begin{center}
    \begin{tikzpicture}[auto,>=latex]
      \node[] (I) {\code{I}};
      \node[below of=I, node distance=.5cm] (O) {\code{O}};
      \node[block, right of=I] (R) {$R$};
      \node[right of=R] (o) {\code{O}};
      \draw[->] (I) -- (R);
      \draw[->] (O) -| (R);
      \draw[->] (R) -- (o);
    \end{tikzpicture} 
    \end{center}
\item Lift this circuit to operate on streams and connect the output to the input in a feedback cycle as follows:

\noindent
\begin{tikzpicture}[auto,>=latex, node distance=.8cm]
  \node[] (Iinput) {\code{I}};
  \node[block, right of=Iinput] (ID) {$\delta_0$};
  \node[block, right of=ID] (II) {$\I$};
  \node[block, right of=II] (f) {$\lift{R}$};
  \node[block, right of=f, node distance=1.5cm] (distinct) {$\lift{\distinct}$};
  \node[block, right of=distinct, node distance=1.5cm] (D) {$\D$};
  \node[block, right of=D] (S) {$\int$};
  \node[right of=S] (output)  {\code{O}};
  \draw[->] (Iinput) -- (ID);
  \draw[->] (ID) -- (II);
  \draw[->] (II) -- (f);
  \draw[->] (f) -- (distinct);
  \draw[->] (distinct) -- node (o) {$o$} (D);
  \draw[->] (D) -- (S);
  \draw[->] (S) -- (output);
  \node[block, below of=distinct, node distance=.7cm] (z) {$\zm$};
  \draw[->] (o) |- (z);
  \draw[->] (z) -| (f);
\end{tikzpicture}

We construct $\lift{R}$ by lifting each operator of the circuit individually 
according to Proposition~\ref{prop:distributivity}.
\end{enumerate}

The inner loop of the circuit computes the fixed point of $R$.  The differentiation
operator $\D$ yields the set of new Datalog facts (changes) computed by each iteration of the loop.
When the set of new facts becomes empty the iterations have completed.
$\int$ computes the value of the fixed point by aggregating these changes.

\begin{theorem}[Recursion correctness]\label{theorem:recursion}
If $\isset(\code{I})$, the output of the circuit above is
the relation $\code{O}$ as defined by the Datalog semantics 
as a function of the input relation \code{I}.
\end{theorem}
\label{proof-recursion}
\begin{proof}
Let us compute the contents of the $o$ stream, produced at the output
of the $\distinct$ operator.  We will show that this stream is composed
of increasing approximations of the value of \code{O}.

We define the following one-argument function: $S(x) = \lambda x . R(\code{I}, x)$.
Notice that the left input of the $\lift{R}$ block is a constant stream
with the value \code{I}.  Due to the stratified nature of the language,
we must have $\ispositive(S)$, so $\forall x . S(x) \geq x$.
Also $\lift{S}$ is time-invariant, so $S(0) = 0$.
From \secref{sec:relational}, the definition of set union we know that
$x \cup y = \distinct(x + y)$.
We get the following system of equations:
$$
\begin{aligned}
o[0] =& S(0) \\
o[t] =& S(o[t-1]) \\
\end{aligned}
$$ 
So, by induction on $t$ we have $o[t] = S^t(0)$, where by 
$S^t$ we mean $\underbrace{S \circ S \circ \ldots \circ S}_{t}$.
$S$ is monotone; thus, if there is a time $k$ such that $S^k(0) = S^{k+1}(0)$, we have 
$\forall j \in \N . S^{k+j}(0) = S^k(0)$.  Applying a derivative to this stream
will then produce a stream that is zero almost everywhere, and integrating
this derivative will return the last distinct value in the stream $o$.

This is essentially the definition of the semantics of a recursive Datalog relation:
$\code{O} = \fix{x}{R(\code{I}, x)}$.
\end{proof}

Note that the use of unbounded data domains (like integers with arithmetic) does 
not guarantee convergence for all programs.

In fact, this circuit implements the standard \defined{na\"{\i}ve evaluation}
algorithm (e.g., see Algorithm~1 in \cite{greco-sldm15}).
Notice that the inner part of the circuit is the incremental
form of another circuit, since it is sandwiched between $\I$ and $\D$ operators.
Using the cycle rule of Proposition~\ref{prop-inc-properties} we can rewrite this circuit as:
%
\begin{equation}
\begin{aligned}
\label{eq:seminaive}
\begin{tikzpicture}[auto,>=latex]
  \node[] (Iinput) {\code{I}};
  \node[block, right of=Iinput] (Idelta) {$\delta_0$};
  \node[block, right of=Idelta] (f) {$\inc{(\lift{R})}$};
  \node[block, right of=f, node distance=1.6cm] (distinct) {$\inc{\lift{(\distinct})}$};
  \node[block, right of=distinct, node distance=1.5cm] (S) {$\int$};
  \node[right of=S] (output)  {\code{O}};
  \node[block, below of=distinct, node distance=.7cm] (z) {$\zm$};
  \draw[->] (Iinput) -- (Idelta);
  \draw[->] (f) -- (distinct);
  \draw[->] (distinct) -- node (o) {} (S);
  \draw[->] (S) -- (output);
  \draw[->] (o) |- (z);
  \draw[->] (z) -| (f);
  \draw[->] (Idelta) -- (f);
\end{tikzpicture}
\end{aligned}
\end{equation}

This last circuit effectively implements the \defined{semi-na\"{\i}ve evaluation}
algorithm (Algorithm~2 from~\cite{greco-sldm15}).  The correctness of semi-na\"{\i}ve
evaluation is an immediate consequence of the cycle rule.
